% ========================================
% PREAMBLE for Probability and Statistics (Classical Grayscale)
% ========================================

% --- Basics ---
\usepackage[utf8]{inputenc}
\usepackage[T1]{fontenc}
\usepackage{textcomp}
\usepackage{graphicx}
\usepackage{float}
\usepackage{booktabs}
\usepackage[shortlabels]{enumitem}
\usepackage{emptypage}
\usepackage{subcaption}
\usepackage{multicol}
\usepackage[usenames,dvipsnames]{xcolor}
\setlength{\headheight}{14.5pt}

% --- Math Packages ---
\usepackage{amsmath, amsfonts, mathtools, amsthm, amssymb}
\usepackage{mathrsfs}
\usepackage{cancel}
\usepackage{bm}
\usepackage[cal=boondoxo]{mathalfa}
\usepackage{systeme}
\usepackage{stmaryrd}

% --- Math Commands ---
\newcommand\N{\mathbb{N}}
\newcommand\R{\mathbb{R}}
\newcommand\Z{\mathbb{Z}}
\renewcommand\O{\emptyset}
\newcommand\Q{\mathbb{Q}}
\newcommand\C{\mathbb{C}}
\newcommand{\samplespace}{\mathcal{S}}
\DeclareMathOperator{\sgn}{sgn}
\let\implies\Rightarrow
\let\impliedby\Leftarrow
\let\iff\Leftrightarrow
\let\epsilon\varepsilon
\newcommand\contra{\scalebox{1.1}{$\lightning$}}
\newcommand\mat[1]{\mathbf{#1}}

% --- Correction Commands ---
\definecolor{correct}{HTML}{555555} % subtle gray
\newcommand\correct[2]{\ensuremath{\:}{\color{red}{#1}}\ensuremath{\to}{\color{black!70}{#2}}\ensuremath{\:}}
\newcommand\green[1]{{\color{correct}{#1}}}

% --- Utilities ---
\usepackage{xifthen}
\newcommand\hr{\noindent\rule[0.5ex]{\linewidth}{0.5pt}}
\newcommand\hide[1]{}

% --- SI Units ---
\usepackage{siunitx}
\sisetup{locale=FR}

% --- TikZ ---
\usepackage{tikz}
\usepackage{tikz-cd}
\usetikzlibrary{intersections, angles, quotes, calc, positioning, arrows.meta, shapes, trees}
\usepackage{pgfplots}
\pgfplotsset{compat=1.13}
\tikzset{
    force/.style={thick, {Circle[length=2pt]}-Stealth, shorten <=-1pt, draw=black!70}
}

% --- Page Layout ---
\usepackage[top=3cm,bottom=3cm,inner=3cm,outer=2.5cm]{geometry}
\setlength{\headheight}{13.6pt}

% --- Theorem Styles ---
\usepackage{thmtools}
\usepackage[framemethod=TikZ]{mdframed}
\mdfsetup{skipabove=1em, skipbelow=0em, innertopmargin=5pt, innerbottommargin=6pt}

\theoremstyle{definition}

% --- Grayscale Theorem Styles ---
\declaretheoremstyle[
    headfont=\bfseries\sffamily,
    bodyfont=\normalfont,
    mdframed={
        linewidth=2pt, rightline=false, topline=false, bottomline=false,
        linecolor=black!50, backgroundcolor=black!5
    }
]{thmgreenbox}

\declaretheoremstyle[
    headfont=\bfseries\sffamily,
    bodyfont=\normalfont,
    mdframed={
        linewidth=2pt, rightline=false, topline=false, bottomline=false,
        linecolor=black!50, backgroundcolor=black!10
    }
]{thmbluebox}

\declaretheoremstyle[
    headfont=\bfseries\sffamily,
    bodyfont=\normalfont,
    mdframed={
        linewidth=2pt, rightline=false, topline=false, bottomline=false,
        linecolor=black!50
    }
]{thmblueline}

\declaretheoremstyle[
    headfont=\bfseries\sffamily,
    bodyfont=\normalfont,
    mdframed={
        linewidth=2pt, rightline=false, topline=false, bottomline=false,
        linecolor=black!50, backgroundcolor=black!10
    }
]{thmredbox}

\declaretheoremstyle[
    headfont=\bfseries\sffamily,
    bodyfont=\normalfont,
    numbered=no,
    mdframed={
        linewidth=2pt, rightline=false, topline=false, bottomline=false,
        linecolor=black!50, backgroundcolor=black!2
    },
    qed=\qedsymbol
]{thmproofbox}

\declaretheoremstyle[
    headfont=\bfseries\sffamily,
    bodyfont=\normalfont,
    numbered=no,
    mdframed={
        linewidth=2pt, rightline=false, topline=false, bottomline=false,
        linecolor=black!50, backgroundcolor=black!2
    },
]{thmexplanationbox}

% --- Theorem Declarations ---
\declaretheorem[style=thmsolutionbox, name=Solution]{solution}
\declaretheorem[style=thmgreenbox, name=Definition]{definition}
\declaretheorem[style=thmbluebox, numbered=no, name=Example]{eg}
\declaretheorem[style=thmbluebox, name=Exercise]{exercise}
\declaretheorem[style=thmredbox, name=Proposition]{prop}
\declaretheorem[style=thmredbox, name=Theorem]{theorem}
\declaretheorem[style=thmredbox, name=Lemma]{lemma}
\declaretheorem[style=thmredbox, numbered=no, name=Corollary]{corollary}
\declaretheorem[style=thmproofbox, name=Proof]{replacementproof}
\renewenvironment{proof}[1][\proofname]{\vspace{-10pt}\begin{replacementproof}}{\end{replacementproof}}
\declaretheorem[style=thmexplanationbox, name=Explanation]{explanation}
\declaretheorem[style=thmblueline, numbered=no, name=Remark]{remark}
\declaretheorem[style=thmblueline, numbered=no, name=Note]{note}

\newtheorem*{notation}{Notation}
\newtheorem*{previouslyseen}{As previously seen}
\newtheorem*{problem}{Problem}
\newtheorem*{observe}{Observe}
\newtheorem*{property}{Property}
\newtheorem*{intuition}{Intuition}

% --- Lecture & Exercise Commands ---
\makeatletter
\newlist{subtasks}{enumerate}{1}
\setlist[subtasks,1]{label=\alph*)}
\def\@lesson{}%
\newcommand{\lesson}[3]{%
    \ifthenelse{\isempty{#3}}{%
        \def\@lesson{Lecture #1}%
    }{%
        \def\@lesson{Lecture #1: #3}%
    }%
    \subsection*{\@lesson}%
    \marginpar{\small\textsf{\mbox{#2}}}%
}
\makeatother

% --- Headers and Footers ---
\usepackage{fancyhdr}
\pagestyle{fancy}
\fancyhead{}
\fancyhead[RO,LE]{MTH 3301}
\fancyfoot[LE,RO]{\thepage}
\fancyfoot[C]{\leftmark}
\renewcommand{\headrulewidth}{0.4pt}

% --- Notes ---
\usepackage{tcolorbox}
\tcbuselibrary{breakable}
\newenvironment{verbetering}{\begin{tcolorbox}[
    arc=0mm, colback=white, colframe=black!50,
    title=Remark, fonttitle=\sffamily, breakable
]}{\end{tcolorbox}}

\newenvironment{noot}[1]{\begin{tcolorbox}[
    arc=0mm, colback=white, colframe=black!30,
    title=#1, fonttitle=\sffamily, breakable
]}{\end{tcolorbox}}

% --- Figure Support ---
\usepackage{import}
\usepackage{pdfpages}
\usepackage{transparent}
\newcommand{\incfig}[1]{%
    \def\svgwidth{\columnwidth}
    \import{./figures/}{#1.pdf_tex}
}
\pdfsuppresswarningpagegroup=1

% --- Hyperlinks ---
\usepackage{hyperref}
\hypersetup{
    colorlinks=false,
    pdfborder={0 0 0}
}
